\section{Discussion}\label{sec:discussion}

\paragraph{The trade.}
Looked at together, the results describe a trade between earning more along the way and finding less at the
end. Sharing raises what the team earns while it looks and lowers what the team eventually finds. A team that
cares about the ceiling seems to do better keeping its agents apart for a while. A team that cares about the
total seems to do better letting them copy each other. The 60-round runs show that this trade is a horizon
effect, and that it resolves in opposite directions depending on the letter. With the letter, it dissolves in
sharing's favour: given time, sharing teams leave the recommended location, reach the same ceiling as agents
working alone, and keep a larger cumulative lead. Without the letter, it dissolves the other way: the sharing
teams' ceiling stays \HAliMaxFull{} points below, and the lone agents' wider search pays back by round 60, so
the cumulative lead of sharing is gone. What sharing buys, on this reading, is a faster climb up whatever hill
the team starts on. Whether that is the right hill is set by where the agents start, and the letter gives every
agent the same start.

\paragraph{The mechanism.}
We think the whole mechanism is that sharing makes the team search less. Once agents can see each other,
they copy whoever is doing best so far. That is good for earnings in the current round. It could be bad for
finding the best location, because the agent doing best so far is usually at a location that is good but
not the best, and the best location pays badly until its recipe is found. In the lemonade framing the
manager's letter seems to make this worse: all three agents start near the letter's location, so the board
shows them three versions of the same wrong answer. Telling agents to stay different keeps exploration at the
working-alone level in both framings, and it brings the max back to the working-alone level in the alien
framing but not in the lemonade framing. One reading is that exploring widely is necessary for finding the
ceiling but not sufficient. One agent at a time explores least of all and still earns the most in the
lemonade framing: each agent sees the two moves before its own in every round and refines them, which is a
good way to climb a hill and a poor way to find out whether there is a higher hill elsewhere.

\paragraph{What is established and what is not.}
Two results are established on this testbed: sharing cuts distinct setups roughly in half, and sharing
raises cumulative earnings. Both hold in every run and both pass the pre-specified test in each framing. The
drop in the max is consistent across runs and framings but is not statistically reliable with eight worlds,
and we report it as a consistent direction, not as a finding. The agents were asked for overall performance
and, in effect, read that as earning well each round; a team told that only its best round counts would have
less reason to copy a known setup, and we would expect a smaller drop under that instruction. No sharing
condition beats working alone on the max; the two that come out level with it do so by one extra world-run
and are within the run-to-run noise.
The paired test over eight worlds has low power for a best-of-sixty measure, so a reader who wants the max
result settled either way should expect it to need several times as many worlds.

\paragraph{Limitations.}\label{sec:limitations}
The study is small: eight worlds for most contrasts, two models, teams of three, a horizon of \NRounds{}
rounds for all but two conditions. The 60-round runs cover working alone and sharing all along only, so how the
periodic board and the other conditions behave in a longer game is not known. The goal line leaves the
objective to the agent. That is the realistic case and it keeps the comparison between conditions clean, but
it means the two measures are read off agents that were optimising neither by instruction, and how far the
drop in the max depends on this is untested; a version of the two main conditions with the objective stated
explicitly, as the total over the game and as the best single round, is the next run. The second model covers
three of
the nine conditions, and on it the drop in the max appears in only one
framing, so the max result should be read as model-dependent until more models are run. Two details of
the alien framing should be known when reading its numbers. First, the shared board is rendered with the
lemonade vocabulary in both framings: its posts name the location and the variables as ``location'',
``sugar'', ``lemon'', ``colour'' and ``price'' and its header describes the agents as stand owners, so an
alien-framing agent that reads the board sees lemonade words there, while its own prompt uses A to E. The
agent's own posts appear on the board, so it can align the two vocabularies, but the alien sharing conditions
are less fully stripped of lemonade priors than the alien working-alone condition. Second, the alien labels
A1 to A3 follow the world's internal ordering of locations rather than a per-episode shuffle, and under that
ordering the first-listed location is never the best in any of the eight worlds; alien agents spend more
attempts at the first-listed location than lemonade agents do. Both details apply equally to every alien
condition, so the within-framing contrasts that this paper reports are unaffected, but we do not compare
levels across framings. Finally, the max is a best-of-three measure. Each agent's own best round is higher
under sharing, so the team-level drop is a statement about the ceiling of the best member, not about the
typical agent.

\paragraph{Next steps.}\label{sec:next}
The first is the objective: working alone and sharing all along with the goal line stated explicitly, as the
total over the game and as the best single round. If the drop in the max survives an instruction to maximise
the best round, the ceiling result becomes a finding rather than a direction. The alien framing changes two things at
once, the labels and the letter, so it cannot say which drives the difference in timing between framings. Two
controls separate them: the lemonade framing without the letter, and the alien framing with the letter. Beyond that, the question of when sharing helps is a question about
the landscape. On a smooth landscape, where local search reaches the top, sharing should help; on a rugged one
it should hurt. Parameterising the worlds by the number of decision variables and their interdependence, as in
NK models \citep{levinthal1997adaptation,lazer2007network}, and running the same conditions at low, medium and
high interdependence with everything else fixed, would turn the present result into a curve. More worlds, more
models, the remaining conditions at the longer horizon, and larger teams are the other obvious extensions.
